\chapter{1915:Richard Willstätter}
\label{chap:chemistry1915}

The Nobel Prize in Chemistry for the year 1915 was awarded to Richard Willstätter.

\textbf{Motivation:}
for his researches on plant pigments, especially chlorophyll

\section{Richard Willstätter}

\subsection*{Early Life and Education}

Richard Willstätter was born on 1872-08-13 in Karlsruhe, Germany.
Richard Martin Willstätter FRS HFRSE was a German organic chemist whose study of the structure of plant pigments, chlorophyll included, won him the 1915 Nobel Prize for Chemistry.

\subsection*{Career and Major Research}

Willstätter was professor at Munich and later director of the Kaiser Wilhelm Institute for Chemistry in Berlin. He determined the constitution of chlorophyll and of the plant pigments, and he studied the anthocyanins, the red and blue pigments of flowers and fruits. He also made important contributions to the chemistry of enzymes.

\subsection*{Nobel Prize-winning Work}

The work for which Richard Willstätter was awarded the Nobel Prize in Chemistry in 1915 was described by the Nobel Committee as: "for his researches on plant pigments, especially chlorophyll".

This research fundamentally advanced the field by making groundbreaking contributions that transformed chemical science.

\subsection*{Significance and Impact}

The impact of Richard Willstätter's work extends far beyond the specific achievement recognized by the Nobel Prize.
These contributions continue to influence chemical research and education today.

\subsection*{Later Career and Honors}

Richard Willstätter continued his scientific work until 1942-08-03, passing away in Locarno, Switzerland.
He received numerous honors including membership in major scientific academies and societies.

\subsection*{Legacy}

The legacy of Richard Willstätter endures through the lasting impact of his scientific contributions.
Richard Willstätter's work continues to be cited and built upon by researchers across the chemical sciences.

\subsection*{Modern Developments}

Willstätter's structural studies of chlorophyll and plant pigments revealed the chemistry of photosynthesis and founded modern pigment research. His work led to the understanding of anthocyanins and carotenoids, which today are studied as dietary antioxidants and natural colorants in foods and cosmetics. Photosynthesis research inspired by Willstätter underpins modern solar-energy and biofuel technologies.

\subsection*{Selected Publications}

Richard Willstätter's major publications include numerous papers in leading journals of the era, detailing the experimental and theoretical work summarized above.

\clearpage

\section*{Chapter Summary}

Richard Willstätter received the prize for his researches on plant pigments, especially chlorophyll.
